\documentclass[11pt]{article}
\usepackage[margin=27mm]{geometry}
\usepackage[T1]{fontenc}
\usepackage{lmodern}
\usepackage{amsmath,amssymb,amsthm,mathtools,mathrsfs}
\usepackage{booktabs,array,microtype}
\usepackage[numbers,sort&compress]{natbib}
\usepackage[colorlinks=true,allcolors=blue]{hyperref}
\usepackage{xr}
\externaldocument[P-]{../main}
\newtheorem{theorem}{Theorem}[section]
\newtheorem{proposition}[theorem]{Proposition}
\newtheorem{lemma}[theorem]{Lemma}
\newtheorem{corollary}[theorem]{Corollary}
\newcommand{\CC}{\mathbb C}
\newcommand{\QQ}{\mathbb Q}
\newcommand{\ZZ}{\mathbb Z}
\newcommand{\FF}{\mathbb F}
\newcommand{\remA}{\operatorname{rem}_{A}}
\newcommand{\len}{\operatorname{length}}
\newcommand{\Xp}{\mathcal X_p}
\newcommand{\dd}{\,\mathrm d}
\numberwithin{equation}{section}
\setlength{\emergencystretch}{2em}
\title{Counts with multiplicity at general orders\\[4pt]
\large Companion note to ``A binomial count of meromorphic travelling waves''}
\input{../authors}
\date{}
\begin{document}
\maketitle

\begin{abstract}
The paper proves its counts when \(2p-1\) is prime, by reduction modulo
that prime. This note contains what is known at the other orders. A
weighted form of B\'ezout's theorem gives the upper bound
\(\binom{2p-1}{p-1}\) at every order, and equality, with pairs counted
with multiplicity, whenever a rigidity condition holds. The same method
gives the counts of pulses, of the pairs of the two pure subfamilies and
of elliptic pairs. The note also proves rigidity when \(2p-1\) is a
prime power and records the computations that establish it for
\(p\le573\), and distinctness of all pairs for \(p\le12\). The text is
that of the first complete version of the paper, where these results
were part of the main line.
\end{abstract}

\section{Notation}

References of the form ``Theorem~P\,5.4'' are to the paper. We use its
notation throughout. The profile equation is
\(P(D)v+v^2/2=0\) with \(P\) monic of degree \(p\)
(equation \eqref{P-eq:main} of the paper). A rational-exponential pair
has the profile \(v_A=s_pA(D)Q\), where \(Q=1/(1+e^{-z})\) and
\(A(\rho)=\rho^d+\alpha_1\rho^{d-1}+\cdots+\alpha_d\), \(d=p-1\). The
discrete convolution is \(S_A=A\star A\), and the pair exists exactly
when \(A\mid S_A\) (Proposition~P\,\ref{P-prop:divisibility}). Writing
\(S_A=AC_A+R_A\), \(R_A=\sum_{r<d}R_r\rho^r\), the matching equations are
\(R_0=\cdots=R_{d-1}=0\), and \(\Xp\) is the affine scheme that they
define. A solution is \emph{simple} if the Jacobian matrix of the
\(R_r\) is invertible there. When \(\Xp\) is finite, its \emph{length}
is the sum of the multiplicities of its points, and \(\Xp\) is
\emph{reduced} when every point is simple, so that the length is the
number of points. We put
\[
 N_p=\binom{2p-1}{p-1},\qquad S_p=\binom{2p-2}{p-2},\qquad
 F_p=\binom{p-1}{\lfloor(p-1)/2\rfloor},\qquad M_p=2S_p .
\]
The variable \(\alpha_j\) has weight \(j\) and \(\rho\) has weight one;
\(H_A(\rho)=\int_0^\rho A(t)A(\rho-t)\dd t\) is the continuous
convolution, and \((\mathrm R_d)\) is the rigidity statement
\eqref{P-eq:rigidity} of the paper: a monic \(A\) of degree \(d\) with
\(A\mid H_A\) is \(\rho^d\). The reflection is
\((\iota A)(\rho)=(-1)^dA(-\rho)\), with the identities of
Proposition~P\,\ref{P-prop:reflection}. For a fixed lattice,
\(\mathcal E_p(g_2,g_3)\) is the elliptic matching system of
Section~P\,\ref{P-sec:elliptic-counting}.

\section{Weighted degree and the count with multiplicities}

B\'ezout's theorem bounds the sum of the multiplicities of the
isolated solutions of \(d\) polynomial equations in \(d\) unknowns
by the product of their degrees. Equality holds when the system has
no solutions at infinity, that is, when the highest-degree parts of
the equations have no common zero other than the origin. The ordinary
degrees of the matching equations give a poor bound. The right
measure of their size is a weighted degree that reflects how the
\(\alpha_j\) scale.
Give \(\alpha_j\) weight \(j\), and \(\rho\) weight one.
The remainder equations have highest weighted degrees
\begin{equation}\label{eq:weights}
 2d+1-r\quad(0\le r<d),\qquad\text{namely }d+2,\ldots,2d+1.
\end{equation}
Their highest weighted parts are the coefficients of
\(\remA H_A\), where
\begin{equation}\label{eq:continuous}
 H_A(\rho)=\int_0^\rho A(t)A(\rho-t)\dd t.
\end{equation}
Indeed, if \(A(\rho)=\sum a_i\rho^i\), then
\[
 H_A(\rho)=\sum_{i,j=0}^{d}
     \frac{i!j!}{(i+j+1)!}a_ia_j\rho^{i+j+1}.
\]
These are precisely the highest weighted terms of the power sums
defining \(S_A\), and monic division preserves the highest weighted
part. Indeed, \(A\) is weighted-homogeneous of weight \(d\), and each
step of long division subtracts a multiple of \(A\) of the same
weight as the term being removed.
In words, the discrete convolution \(S_A\) is a lower-order
perturbation of the continuous convolution \(H_A\), and the behaviour
of the matching equations at infinity is governed by the continuous
problem \(A\mid H_A\).

For this degree, denote by \((\mathrm{R}_d)\) the condition
\begin{equation}\label{eq:rigidity}
 A\text{ monic},\quad\deg A=d,\quad A\mid H_A
 \quad\Longrightarrow\quad A(\rho)=\rho^d.
\end{equation}
It says that the continuous problem has only the trivial solution.
By the preceding paragraph, it is exactly the condition that the
matching equations have no solution at weighted infinity.

\begin{theorem}[Count with algebraic multiplicities]\label{thm:length}
For every \(p\ge2\), the sum of multiplicities of the isolated points
of \(\Xp\) is at most
\[
 N_p=\binom{2p-1}{p-1}.
\]
If \((\mathrm R_{p-1})\) holds, then \(\Xp\) is finite and
\begin{equation}\label{eq:length}
 \len\Xp=N_p.
\end{equation}
This equality does not require simplicity.
\end{theorem}

\begin{proof}
Weighted degree is turned into ordinary degree by the finite
substitution
\begin{equation}\label{eq:weighted-cover}
 \alpha_j=y_j^j,\qquad 1\le j\le d.
\end{equation}
A monomial of weight \(w\) in the \(\alpha_j\) becomes a monomial of
ordinary degree \(w\) in the \(y_j\), so the pulled-back equations
have ordinary degrees at most \(d+2,\ldots,2d+1\), and their
highest-degree parts are the pull-backs of the highest weighted parts.
B\'ezout's inequality for isolated solutions bounds the sum of their
multiplicities by
\((d+2)\cdots(2d+1)\); see \citet[Chapter~VIII]{Mondal2021}.
If \((\mathrm R_d)\) holds, the leading homogeneous forms have
no nonzero common zero. The homogenized system has no point at
infinity, and ordinary B\'ezout gives equality
\citep{CoxLittleOShea2005}.

It remains to return from the \(y_j\) to the \(\alpha_j\).
A solution in the \(\alpha_j\) with all \(\alpha_j\ne0\) has
\(1\cdot2\cdots d=d!\) preimages in the \(y_j\), one for each
choice of the roots \(y_j=\alpha_j^{1/j}\), each with the original
multiplicity. In general, for each solution the sum of the
multiplicities of its preimages is \(d!\) times its original
multiplicity, including at branch points where some \(\alpha_j\)
vanish. Thus the substitution multiplies the total length by \(d!\).
Algebraically, \(\CC[y_1,\ldots,y_d]\) is a free module over
\(\CC[\alpha_1,\ldots,\alpha_d]\), with basis
\[
 y_1^{e_1}\cdots y_d^{e_d},\qquad 0\le e_j<j.
\]
This remains true after imposing the equations and restricting to
an isolated component.
Dividing the B\'ezout count by \(d!\) gives
\[
 \frac{(d+2)\cdots(2d+1)}{d!}=\binom{2d+1}{d}.
\]
The no-infinity condition also gives finiteness of the whole affine
scheme, completing the proof.
\end{proof}

For \(p=3\), the variable weights are \(1,2\) and the equation
weights are \(4,5\), so the calculation is
\((4\cdot5)/(1\cdot2)=10\).


\section{Pulses}

The endpoints of \(v_A\) are \(0\) and \(s_pA(0)\); the pairs with
\(A(0)=0\) are the pulses.

\begin{proposition}[Equal endpoints]\label{cor:endpoints}
Assume \((\mathrm R_{p-1})\), and let
\(\mathcal Z_p=\Xp\cap\{A(0)=0\}\), with its induced scheme structure.
Then
\begin{equation}\label{eq:pulse-length}
 \len\mathcal Z_p=S_p:=\binom{2p-2}{p-2}.
\end{equation}
If \(\Xp\) is reduced, these are \(S_p\) distinct pairs and
\(P_A(0)\ne0\) at each of them.
In particular, under the prime-index hypothesis, exactly
\[
 \binom{2p-2}{p-2}
\]
normalized pairs have equal zero endpoints. The remaining
\(\binom{2p-2}{p-1}\) have unequal endpoints.
These are counts over \(\CC\).
\end{proposition}

\begin{proof}
Polynomial summation gives \(S_A(0)=-A(0)^2\), so the
constant remainder satisfies the polynomial identity
\begin{equation}\label{eq:constant-remainder}
 R_0=-A(0)\bigl(A(0)+C_A(0)\bigr).
\end{equation}
Set \(\alpha_d=A(0)=0\).
The equation \(R_0=0\) disappears, leaving \(d-1\) equations
of weights \(d+2,\ldots,2d\) in variables of weights
\(1,\ldots,d-1\).
Their leading forms have no nonzero common zero by
\((\mathrm R_d)\). The weighted calculation in
Theorem~\ref{thm:length} gives
\[
 \len\mathcal Z_p
 =\frac{(d+2)\cdots(2d)}{(d-1)!}
 =\binom{2d}{d-1}.
\]
For \(d=1\) the system is empty and its length is one.

If \(\Xp\) is finite and reduced, its coordinate algebra is a product
of copies of \(\CC\); the same is true of any quotient, in particular
that of \(\mathcal Z_p\), so \(\mathcal Z_p\) is reduced as well.
At a point with \(A(0)=C_A(0)=0\), the differential of
\eqref{eq:constant-remainder} vanishes.
Such a point cannot be simple in the square matching system.
Therefore \(C_A(0)\ne0\), and
\(P_A(0)=s_pC_A(0)/2\ne0\), at every point of \(\mathcal Z_p\)
when \(\Xp\) is reduced.
The prime-index assertions now follow from Theorem~P\,\ref{P-thm:prime}
and \(N_p-S_p=\binom{2p-2}{p-1}\).
\end{proof}


\section{The two pure subfamilies}

Let
\[
 \Xp^{\mathrm{sym}}=\Xp\cap\{\alpha_j=0\text{ for odd }j\}
\]
be the fixed locus of \(\iota\), with its induced scheme structure, and
put
\begin{equation}\label{eq:Fp}
 F_p=\binom{p-1}{\lfloor (p-1)/2\rfloor}.
\end{equation}
Its first values are \(1,2,3,6,10,20\) for \(p=2,\ldots,7\);
Table~P\,\ref{P-tab:named} lists the equations to which they belong.

\begin{theorem}[Count in the two pure subfamilies]\label{thm:symmetric}
\begin{enumerate}
\item[\textup{(i)}] A point \(A\) of \(\Xp\) lies in
\(\Xp^{\mathrm{sym}}\) if and only if \(P_A\) is reflection-symmetric.
Thus \(\Xp^{\mathrm{sym}}\) consists of the rational-exponential pairs
of the purely dispersive equations when \(p\) is even, and of the
purely dissipative equations when \(p\) is odd.
\item[\textup{(ii)}] If \((\mathrm R_{p-1})\) holds, then
\(\len\Xp^{\mathrm{sym}}=F_p\).
\item[\textup{(iii)}] If, in addition, \(\Xp\) is reduced, then
\(\Xp^{\mathrm{sym}}\) consists of \(F_p\) distinct pairs, each a simple
solution of the matching equations of its subfamily. If \(2p-1\) is
prime, they are the points whose label in (P\,\ref{P-eq:prime-labels})
satisfies \(S=-S\).
\item[\textup{(iv)}] For even \(p\), every pair of
\(\Xp^{\mathrm{sym}}\) is a pulse. For odd \(p\), every pair of
\(\Xp^{\mathrm{sym}}\) is a front if \(\Xp\) is reduced, and every pair
with real coefficients is a front without any hypothesis.
\end{enumerate}
\end{theorem}

In plain terms, the pairs that coincide with their mirror image are
exactly the pairs of the purely dispersive equations, for even \(p\),
or of the purely dissipative equations, for odd \(p\). Counted with
multiplicity there are \(F_p\) of them whenever rigidity holds, in
particular for \(p\le573\). They are \(F_p\) distinct pairs whenever all
\(N_p\) pairs for that \(p\) are distinct, in particular when \(2p-1\)
is prime. The rational-exponential waves of a purely dispersive
equation are pulses, for every even \(p\). Those of a purely
dissipative equation are fronts; this is proved whenever all \(N_p\)
pairs are distinct, and for waves with real coefficients at every odd
\(p\). Elliptic waves occur in both subfamilies; the distinction
between pulses and fronts does not apply to them.

\begin{proof}
\emph{(i).} If \(\iota A=A\), then (P\,\ref{P-eq:reflection-pair}) gives
\(P_A(\rho)=(-1)^p(P_A(-\rho)+s_pA(0))\). For even \(p\) the
polynomial \(A\) is odd, so \(A(0)=0\) and \(P_A\) is even. For odd
\(p\) the sum \(P_A(\rho)+P_A(-\rho)\) is the constant \(-s_pA(0)\).
In both cases \(P_A\) is reflection-symmetric, and
\begin{equation}\label{eq:symmetric-endpoint}
 s_pA(0)=-2P_A(0)\qquad(p\text{ odd}).
\end{equation}

Conversely, let \(P=P_A\) be reflection-symmetric, and apply
Corollary~P\,\ref{P-cor:symmetric-waves} at the pole \(i\pi\), using the
period \(2\pi i\). If \(p\) is even, then \(v_A(-z)=v_A(z)\). Letting
\(z\to-\infty\) gives \(A(0)=0\), and (P\,\ref{P-eq:reflection-pair})
becomes \(v_{\iota A}=v_A\). If \(p\) is odd, then
\(-v_A(-z)=v_A(z)+2P(0)\), and (P\,\ref{P-eq:reflection-pair}) becomes
\(v_{\iota A}(z)=v_A(z)+2P(0)+s_pA(0)\). Letting \(z\to-\infty\) shows
that the constant vanishes. In both cases \(v_{\iota A}=v_A\), hence
\(\iota A=A\) by the uniqueness of the form (P\,\ref{P-eq:normalized-A}).

\emph{(ii).} On the subspace \(\{\alpha_j=0,\ j\text{ odd}\}\) we have
\(\iota A=A\), and (P\,\ref{P-eq:reflection-S}) shows that \(R_A\) is an odd
polynomial. The equations \(R_r=0\) with \(r\) even therefore hold
identically there. With \(m=\lfloor d/2\rfloor\), there remain the
\(m\) equations \(R_1,R_3,\ldots,R_{2m-1}\) in the \(m\) unknowns
\(\alpha_2,\alpha_4,\ldots,\alpha_{2m}\). By \eqref{eq:weights} their
weights are
\[
 \begin{array}{c|c}
 \text{variables}&2,4,\ldots,2m\\
 \text{equations}&2d,2d-2,\ldots,2(d-m+1).
 \end{array}
\]
Their highest weighted parts are the restrictions of those of the full
system. The highest weighted parts of the omitted equations vanish
identically on the subspace, because \(H_A\) is an odd polynomial when
\(\iota A=A\). A common zero of the restricted leading
forms is therefore a polynomial \(A\) with \(A\mid H_A\), and
\((\mathrm R_d)\) gives \(A=\rho^d\). The finite-cover argument of
Theorem~\ref{thm:length}, with \(\alpha_{2i}=y_i^{2i}\), gives
\[
 \len\Xp^{\mathrm{sym}}
 =\prod_{i=1}^{m}\frac{2(d-i+1)}{2i}=\binom dm .
\]

\emph{(iii).} A closed subscheme of a finite reduced scheme is reduced,
as in the proof of Proposition~\ref{cor:endpoints}, and a reduced point
of a square system is a simple solution of it. The map
(P\,\ref{P-eq:iota}) negates the roots of \(\overline A\), so it acts on the
labels by \(S\mapsto-S\), and the labelling is a bijection.

\emph{(iv).} For even \(p\) the polynomial \(A\) is odd, so \(A(0)=0\).
Let \(p\) be odd, and suppose that a point of \(\Xp^{\mathrm{sym}}\) is
a pulse, \(A(0)=0\). Then \(P_A(0)=0\) by
\eqref{eq:symmetric-endpoint}. If \(\Xp\) is reduced, this contradicts
Proposition~\ref{cor:endpoints}. If \(A\) has real coefficients, then
\(v_A\) is real on the real axis, and it tends to zero with all its
derivatives at both ends. Integrating (P\,\ref{P-eq:main}) along the real
axis, and using \(P_A(0)=0\), gives \(\int v_A^2\dd z=0\). Hence
\(v_A=0\), which is impossible.
\end{proof}


\section{Elliptic pairs}

\begin{theorem}[Fixed-lattice elliptic enumeration]\label{thm:elliptic-count}
Assume \((\mathrm R_{p-1})\).
For every fixed nonsingular lattice,
\begin{equation}\label{eq:elliptic-length}
 \len\mathcal E_p(g_2,g_3)
   =M_p:=2\binom{2p-2}{p-2}.
\end{equation}
If, in addition, \(\Xp\) is reduced, then
\(\mathcal E_p(g_2,g_3)\) is reduced for
\((g_2,g_3)\) in a nonempty Zariski-open subset of
\(\{g_2^3-27g_3^2\ne0\}\).
In particular, if \(2p-1\) is prime, a generic fixed lattice
admits exactly \(M_p\) normalized elliptic operator--profile
pairs, all simple.
\end{theorem}

Thus the elliptic assertion on an arbitrary lattice is a count
with multiplicities. Exceptional lattices, on which distinct pairs
merge, do occur and cannot be discarded when counting distinct
solutions; Section~P\,\ref{P-sec:ks} determines them at order three.

The proof has five steps. We eliminate the operator, which leaves
\(p-1\) equations in the \(p-1\) profile unknowns. We assign weights to the unknowns, under which the highest weighted
parts of the equations are obtained by setting \(g_2=g_3=0\). We show
that at this cusp the equations have only the trivial common zero;
this is where \((\mathrm R_{p-1})\) enters, and it means that there
are no solutions at weighted infinity for any \(g_2,g_3\). Weighted B\'ezout then gives the length \(M_p\) on every
lattice. Finally, at a nodal cubic the elliptic problem becomes the
equal-endpoint case of the rational-exponential problem, whose
solutions are simple by Proposition~\ref{cor:endpoints} when \(\Xp\)
is reduced, and simplicity spreads from the nodal parameter value to
nearby nonsingular parameters.

\begin{proof}
\emph{Step 1: eliminating the operator.}
The coefficient at pole order \(2p\) vanishes by the choice of
\(s_p\).
The coefficients at orders \(2p-1,\ldots,p\) determine
\(a_{p-1},\ldots,a_0\) successively.
At each step the coefficient of the new \(a_j\) is the nonzero
rational leading coefficient of \(D^jV\).
This is a polynomial elimination, valid over any
\(\QQ\)-algebra.
The residual has pole order at most \(p-1\) and zero residue.
It lies in the span of
\[
 \{1\}\ \cup\ \{\mathscr D^jX:0\le j\le p-3\}.
\]
Hence \(p-1\) equations remain in the \(p-1\) profile
coefficients \(u_j,b\).

\emph{Step 2: weights.}
Give \(D,X,g_2,g_3\) weights \(1,2,4,6\), respectively, and give
the profile weight \(p\).
Then \(u_j\) has weight \(p-j-2\), and \(b\) has weight \(p\).
After the triangular elimination, the residual principal-part
coefficients and the residual Laurent constant have weights
\begin{equation}\label{eq:elliptic-weights}
 \begin{array}{c|c}
 \text{variables}&1,2,\ldots,p-2,p\\
 \text{equations}&p+1,p+2,\ldots,2p-2,2p.
 \end{array}
\end{equation}
The Laurent expansions and the eliminations are polynomial in
\(g_2,g_3\), so these statements extend to the singular cubics.
For fixed \(g_2,g_3\), the highest weighted parts of the equations in
the profile variables are obtained by setting \(g_2=g_3=0\), since
\(g_2\) and \(g_3\) carry positive weight.

\emph{Step 3: the cusp.}
At the cusp \(g_2=g_3=0\), \(X=z^{-2}\) and \(Y=-2z^{-3}\), so the
profile is a rational function.
Write \(V=b+f(z)\), where \(f\) has its sole pole at zero
and vanishes at infinity.
The constant and decaying parts of the equation give
\[
 a_0b+\frac12b^2=0,\qquad
 (P+b)(D)f+\frac12f^2=0.
\]
The leading term of \(f\) at infinity forces \(a_0+b=0\).
It follows that \(b=a_0=0\).
Under the polynomial Laplace correspondence
\[
 \sum_k A_k\rho^k\ \longmapsto\
       \sum_k k!A_k z^{-k-1},
\]
differentiation corresponds to multiplication by \(-\rho\),
and multiplication of the rational functions corresponds to
continuous convolution.
After a nonzero scalar normalization, the equation for \(f\)
therefore gives \(A\mid H_A\) for a monic \(A\) of degree \(p-1\).
By \((\mathrm R_{p-1})\), \(A=\rho^{p-1}\).
Thus \(f=c_\ast z^{-p}\), and all the variables in
(P\,\ref{P-eq:elliptic-enumeration-profile}) vanish.
The leading equations have no nonzero common zero.

\emph{Step 4: the length.}
The same finite-cover B\'ezout argument as in
Theorem~\ref{thm:length}, now with the weights
\eqref{eq:elliptic-weights}, gives
\[
 \len\mathcal E_p(g_2,g_3)
 =\frac{(p+1)(p+2)\cdots(2p-2)(2p)}{(p-2)!\,p}
 =2\binom{2p-2}{p-2}.
\]
For \(p=3\), this is \((4\cdot6)/(1\cdot3)=8\).
Empty products are one, so the calculation includes \(p=2\).
It proves this length at every \((g_2,g_3)\), including the
singular parameters used below; at the cusp itself the whole length
sits at the single solution where every unknown vanishes.
It also makes the whole coefficient family finite over
\(\CC[g_2,g_3]\): sufficiently large weighted monomials reduce
to bounded weight using the leading equations, which are
independent of \(g_2,g_3\). This finiteness is what allows
simplicity to be transferred from the nodal parameter value to nearby
nonsingular parameters in the last step.

\emph{Step 5: the node.}
To prove generic simplicity, consider the nodal parameters
\begin{equation}\label{eq:nodal-parameters}
 g_2=\frac1{12},\qquad g_3=-\frac1{216}.
\end{equation}
In the logistic coordinate they are realized by
\[
 X=Q^2-Q+\frac1{12},\qquad
 Y=(2Q-1)Q(1-Q),\qquad DQ=Q(1-Q).
\]
The nodal coordinate ring is
\[
 \CC[X,Y]=\{f\in\CC[Q]:f(0)=f(1)\}.
\]
Indeed the right side is
\(\CC+Q(1-Q)\CC[Q]\); it is generated by \(X,Y\), since,
with \(S=Q(1-Q)\),
\(S=1/12-X\), \(SQ=(Y+S)/2\), and \(Q^2=Q-S\).
The pole filtration is the ordinary polynomial degree in \(Q\).
Consequently the nodal specialization of
(P\,\ref{P-eq:elliptic-enumeration-profile}) has the unique form
\[
 V=\beta+v_A,\qquad A\text{ monic of degree }p-1,\quad A(0)=0,
\]
where \(\beta\) is the common endpoint.
The coefficient change is invertible and affine over \(\QQ\).

The matching equations become, as equations of schemes,
\begin{equation}\label{eq:nodal-matching}
 A\in\mathcal Z_p,\qquad
 P=P_A-\beta,\qquad
 \beta\bigl(\beta-2P_A(0)\bigr)=0.
\end{equation}
To see this directly, the constant equation is
\(P(0)\beta+\beta^2/2=0\).
The remaining polynomial identity is
\((P+\beta)(D)v_A+v_A^2/2=0\), equivalent by
Proposition~P\,\ref{P-prop:divisibility} to \(A\mid S_A\) and
\(P+\beta=P_A\).
The coefficient changes and monic division used here work over
arbitrary \(\CC\)-algebras, so the correspondence preserves
local multiplicities.

If \(\Xp\) is reduced, Proposition~\ref{cor:endpoints} gives
\(S_p\) simple points of \(\mathcal Z_p\), with
\(P_A(0)\ne0\) at every point.
Each quadratic in \eqref{eq:nodal-matching} then has two
distinct roots. The nodal fiber has \(2S_p=M_p\) simple points.
In the finite coefficient family, the image of the closed locus
where the matching Jacobian vanishes is closed.
Its complement contains \eqref{eq:nodal-parameters}, and hence
meets the nonsingular-lattice locus in a nonempty open set.
Every fiber in that open set is reduced, proving the theorem.
\end{proof}

The factor two in \(M_p\) retains the two choices of background in
(P\,\ref{P-eq:shift}). That transformation preserves the evolution equation
and the lattice; it changes the operator \(P\) and the speed. At the
nodal fiber it exchanges the two roots in
\eqref{eq:nodal-matching}, which are distinct when \(\Xp\) is reduced.
In that case the two backgrounds are distinct on a generic lattice, and
the \(M_p\) pairs represent \(M_p/2\) elliptic waves, up to
translations and Galilean shifts, each counted with both of its
backgrounds. On special nonsingular lattices the two backgrounds may
coalesce, as Section~P\,\ref{P-sec:ks} shows.

\begin{corollary}\label{cor:elliptic-symmetric}
Let \(p\) be even and assume \((\mathrm R_{p-1})\). On every fixed
nonsingular lattice, the elliptic pairs whose operator is an even
polynomial form a closed subscheme
\(\mathcal E_p^{\mathrm{sym}}(g_2,g_3)\) of
\(\mathcal E_p(g_2,g_3)\), defined by \(u_j=0\) for odd \(j\), and
\begin{equation}\label{eq:elliptic-symmetric-length}
 \len\mathcal E_p^{\mathrm{sym}}(g_2,g_3)=2F_p .
\end{equation}
If \(\Xp\) is reduced, \(\mathcal E_p^{\mathrm{sym}}(g_2,g_3)\) is
reduced for \((g_2,g_3)\) in a nonempty Zariski-open set of nonsingular
parameters.
\end{corollary}

\begin{proof}
An even operator has an even profile, by
Corollary~P\,\ref{P-cor:symmetric-waves}. Conversely, if the profile
(P\,\ref{P-eq:elliptic-enumeration-profile}) is even, the odd part of
(P\,\ref{P-eq:main}) is \(\sum_{j\text{ odd}}a_jD^jV=0\); the functions
\(D^jV\) have distinct pole orders, so \(a_j=0\) for odd \(j\).
On the subspace \(u_j=0\), \(j\) odd, the residual of Step~1 is even,
and there remain the equations and unknowns of even weight in
\eqref{eq:elliptic-weights}:
\[
 \begin{array}{c|c}
 \text{variables}&2,4,\ldots,p-2,p\\
 \text{equations}&p+2,p+4,\ldots,2p-2,2p.
 \end{array}
\]
The omitted leading forms vanish identically on the subspace, so the
remaining ones have only the trivial common zero, by Step~3.
Weighted B\'ezout gives
\[
 \frac{(p+2)(p+4)\cdots(2p-2)\,(2p)}{2\cdot4\cdots(p-2)\,p}
 =2\binom{p-1}{p/2-1}=2F_p .
\]
At the node, reflection is \(Q\mapsto1-Q\). An even nodal profile is
\(\beta+v_A\) with \(\iota A=A\), so the nodal fiber is
\eqref{eq:nodal-matching} with \(A\in\Xp^{\mathrm{sym}}\). If \(\Xp\)
is reduced, it has \(2F_p\) simple points, by
Theorem~\ref{thm:symmetric} and Proposition~\ref{cor:endpoints}, and
Step~5 applies without change.
\end{proof}


\section{What is known at each order}

Each formula needs two facts at a given \(p\). The first is the
rigidity condition \((\mathrm R_{p-1})\) of Section~P\,\ref{P-sec:counting}.
It gives \(N_p\), \(M_p\), \(F_p\) and \(2F_p\) as counts with
multiplicity: the multiplicities of the pairs add up to these numbers.
The second is that every rational-exponential pair is simple, that is,
\(\Xp\) is reduced. Then the pairs are distinct, and the elliptic pairs
are distinct on a generic lattice. On a special lattice several
elliptic pairs can merge, while the count with multiplicity stays the
same; Section~P\,\ref{P-sec:ks} shows this at order three.
Table~\ref{tab:status} records what is known about each fact.

\begin{table}[htbp]
\centering
\caption{What is known. For elliptic pairs, distinctness is asserted
only on a generic lattice.}
\label{tab:status}
\begin{tabular}{@{}p{0.36\linewidth}p{0.58\linewidth}@{}}
\toprule
Statement&Known for\\
\midrule
Counts with multiplicity&\(2p-1\) a prime or a prime power (proved);
every \(p\le73\) (exact certificate); every \(p\le573\) (finite-field
computations)\\[2pt]
All \(N_p\) pairs distinct&\(2p-1\) prime (proved); \(p=5,8,11\) (exact
computations); hence every \(p\le12\)\\[2pt]
All \(F_p\) pairs of the pure subfamily distinct&whenever all \(N_p\)
pairs are distinct; every \(p\le14\) (exact computations)\\[2pt]
Both statements, every \(p\)&conjectured\\
\bottomrule
\end{tabular}
\end{table}

The prime-index argument excluded solutions at infinity by reducing
the continuous convolution modulo \(\ell\). The same argument proves
\eqref{eq:length} whenever \(2p-1=\ell^e\) is an odd prime power
(Appendix~P\,\ref{P-app:verification}). It gives the count with
multiplicity, but it does not show that the pairs are distinct.

More generally, the exact arithmetic certificates supplied with the
paper establish \((\mathrm R_d)\) for every \(d\le72\), hence
\begin{equation}\label{eq:coverage}
 \len\Xp=\binom{2p-1}{p-1},\qquad
 \len\mathcal E_p(g_2,g_3)=2\binom{2p-2}{p-2}
 \qquad(2\le p\le73).
\end{equation}
In the same range, Theorem~\ref{thm:symmetric} and
Corollary~\ref{cor:elliptic-symmetric} give
\(\len\Xp^{\mathrm{sym}}=F_p\) and, for even \(p\),
\(\len\mathcal E_p^{\mathrm{sym}}(g_2,g_3)=2F_p\).
Together with the ancillary proof, the certificates also cover infinite
families of higher orders.
Finite-field Gr\"obner computations, checked by independent pipelines,
establish \((\mathrm R_d)\) for every \(d\le572\); they extend
\eqref{eq:coverage} and the two symmetric lengths, as counts with
multiplicity, to every \(2\le p\le573\) (repository, directory
\texttt{certificates/rigidity\_572/}).
The certificates hold for all complex coefficients, not only real or
rational ones. The ancillary proof and the full certificate are described
in Appendix~P\,\ref{P-app:verification}.

Equation \eqref{eq:coverage} is a statement about counts with
multiplicity only. At composite indices, exact computations show that
all pairs are simple at \(p=5\), \(p=8\) and \(p=11\)
(Appendix~P\,\ref{P-app:verification}). With Theorem~P\,\ref{P-thm:prime}, all
pairs are therefore distinct for every \(p\le12\), and
Theorem~\ref{thm:elliptic-count} gives \(M_p\) distinct elliptic pairs
on a generic lattice at these orders. Inside the two pure subfamilies
the computation is much smaller, and it shows that the \(F_p\) pairs
are distinct for every \(p\le14\), which includes the composite indices
\(9\), \(15\), \(21\), \(25\) and \(27\).

\section{Reduction to characteristic \(\ell\), and prime powers}

The lemma transfers the absence of solutions at infinity from
characteristic \(\ell\) to characteristic zero.
Let homogeneous polynomials \(F_1,\ldots,F_m\), with coefficients
in \(\ZZ_{(\ell)}\), have no nonzero common zero over
\(\overline{\FF}_\ell\).
Then they have no nonzero common zero over \(\CC\).
Indeed a complex projective zero yields a zero over
\(\overline{\QQ}\) after fixing one nonzero coordinate.
Choose a place above \(\ell\), extend it if necessary, and scale
all coordinates by the inverse of a coordinate of smallest valuation.
The coordinates are integral and one is a unit.
Their reductions give a nonzero common zero over
\(\overline{\FF}_\ell\), contradicting the hypothesis.
The argument needs the algebraic closure of the residue field;
a search only over \(\FF_\ell\) would not suffice.

\paragraph{Prime powers.}
Let \(2p-1=\ell^e\) be an odd prime power. For
\(\deg A=(\ell^e-1)/2\), the normalized continuous convolution
satisfies
\[
 \ell^eH_A(\rho)\equiv u\rho^{\ell^e}\pmod\ell,\qquad u\ne0.
\]
Every bilinear coefficient other than the leading one has
denominator valuation at most \(e-1\), whereas the leading
coefficient has valuation \(-e\).
This follows either from the factorial formula or from
\[
 \frac{i!j!}{(i+j+1)!}
   =\sum_{r=0}^{j}\frac{(-1)^r\binom jr}{i+r+1}.
\]
The lemma above then gives \((\mathrm R_{p-1})\), as in Step~3 of the
proof of Theorem~P\,\ref{P-thm:prime}. This argument controls only the
leading system, so it gives no simple-root lifting statement for the
discrete matching system.


The computations referred to above are described in
Appendix~P\,\ref{P-app:verification} of the paper and in the repository
\url{https://github.com/migita/binomial-count-meromorphic-waves},
directories \texttt{certificates/} and \texttt{paper/anc/}.

{\small\raggedright
\bibliographystyle{plainnat}
\bibliography{../references}
}
\end{document}
